\documentclass[../main.tex]{subfiles}
\begin{document}

\section{引言}

在本章中，我们定义了非有理B样条曲线和曲面，研究它们的性质，并推导出它们的导数表达式。为简洁起见，本章其余部分将省略“非有理”一词。主要目标是获得对B样条曲线和曲面的直观理解，为此，读者应仔细研究本章中给出的许多示例和图形。我们还提供了计算B样条曲线和曲面上的点及其导数的算法。Gordon和Riesenfeld首次提出使用B样条定义计算机辅助几何设计的曲线和曲面 \cite{Gord74b,Ries73}。目前，许多关于曲线和曲面的书籍中都涵盖了B样条技术 - 参见 \cite{DeBo78,Mort85,Bart87,Fari93,Yama88,Hosc93,Su89,Roge90,Beac91}。



\section{B样条曲线的定义和性质}\BookIndexMark{p0053}

p次B样条曲线定义为

\[
\mathbf{C}\left( u\right)  = \mathop{\sum }\limits_{{i = 0}}^{n}{N}_{i,p}\left( u\right) {\mathbf{P}}_{i}\;a \leq  u \leq  b \tag{3.1}
\]

其中 \(\left\{  {\mathbf{P}}_{i}\right\}\) 是控制点，而 \(\left\{  {{N}_{i,p}\left( u\right) }\right\}\) 是在非周期(且非均匀)节点向量上定义的 \(p\) 次B样条基函数(公式 [2.5])。

\[
U = \left\{  {\underset{p + 1}{\underbrace{a,\ldots ,a}},{u}_{p + 1},\ldots ,{u}_{m - p - 1},\underset{p + 1}{\underbrace{b,\ldots ,b}}}\right\}
\]

\(\left( {m + 1\text{ knots }}\right)\) 。除非另有说明，我们假设 \(a = 0\) 和 \(b = 1\) 。由 \(\left\{  {\mathbf{P}}_{i}\right\}\) 形成的多边形称为控制多边形\BookIndexMark{p0054}。B样条曲线的示例(在某些情况下与其相应的基函数一起)显示在图3.1-3.14中。

计算固定 \(u\) 值下的B样条曲线上的点需要三个步骤:

1. 找到 \(u\) 所在的节点区间(算法A2.1)；

2. 计算非零基函数(算法A2.2)；

3. 将非零基函数的值与相应的控制点相乘。

考虑第2.5节的例子Ex2.3，具有 \(U = \{ 0,0,0,1,2,3,4,4,5,5,5\}\) 、 \(u = 5/2\) 和 \(p = 2\) 。那么 \(u \in  \left\lbrack  {{u}_{4},{u}_{5}}\right)\) ，并且

\[
{N}_{2,2}\left( \frac{5}{2}\right)  = \frac{1}{8}\;{N}_{3,2}\left( \frac{5}{2}\right)  = \frac{6}{8}\;{N}_{4,2}\left( \frac{5}{2}\right)  = \frac{1}{8}
\]

与控制点相乘得到

\[
\mathbf{C}\left( \frac{5}{2}\right)  = \frac{1}{8}{\mathbf{P}}_{2} + \frac{6}{8}{\mathbf{P}}_{3} + \frac{1}{8}{\mathbf{P}}_{4}
\]

算法如下。

% book-code: algorithm A3.1
\par\medskip\noindent\textbf{算法 A3.1（伪代码）}\BookIndexMark{p0055}\par
\begin{minted}[escapeinside=@@,mathescape]{c}
CurvePoint(n, p, U, P, u, C)
{ /* Compute curve point */
    /* Input: n, p, U, P, u */
    /* Output: C */
    span = FindSpan( n,p,u,U) ;
    BasisFuns(span, u, p, U, N);
    C = 0.0 ;
    for (i = 0; i<=p; i++)
        C = C + N[i]*P[span-p+i];
}
\end{minted}

现在我们列出B样条曲线的一些性质。这些性质源于第2章中给出的关于函数 \({N}_{i,p}\left( u\right)\) 的性质。令 \(\mathbf{C}\left( u\right)\) 由方程(3.1)定义。

P3.1 如果 \(n = p\) 和 \(U = \{ 0,\ldots ,0,1,\ldots ,1\}\) ，那么 \(\mathbf{C}\left( u\right)\) 是一条贝塞尔曲线(图3.1)；

\begin{center}
\begin{minipage}{14cm}
\vspace*{.6\baselineskip}
\centering
\BookFigure[max width=\linewidth]{ch03/fig_3_01.tex}
\par\smallskip
图 3.1. 在 \(U = \{ 0,0,0,0,1,1,1,1\}\) 上的三次B样条曲线，即三次贝塞尔曲线。
\end{minipage}
\end{center}

P3.2 \(\mathbf{C}\left( u\right)\) 是一条分段多项式曲线(因为 \({N}_{i,p}\left( u\right)\) 是分段多项式)；度数 \(p\) 、控制点数量 \(n + 1\) 和节点数量 \(m + 1\) 之间的关系为

\[
m = n + p + 1 \tag{3.2}
\]



(见第2.4节)。图3.2和3.3显示了与各个节点区间对应的B样条曲线的基函数和截面；在这两幅图中，交替的实线/虚线段对应于定义曲线的不同多项式(节点区间)。

\begin{center}
\begin{minipage}{18cm}
\vspace*{.6\baselineskip}
\centering
\BookFigure[max width=\linewidth]{ch03/fig_3_02.tex}
\par\smallskip
图 3.2. (a) 三次基函数 \(U = \{ 0,0,0,0,1/4,1/2,3/4,1,1,1,1\}\) ; (b) 使用图 3.2a 的基函数构造的三次曲线。
\end{minipage}
\end{center}

\begin{center}
\begin{minipage}{18cm}
\vspace*{.6\baselineskip}
\centering
\BookFigure[max width=\linewidth]{ch03/fig_3_03.tex}
\par\smallskip
图 3.3. (a) 在 \(U = \{ 0,0,0,1/5,2/5,3/5,4/5,1,1,1\}\) 上的二次基函数；(b) 使用图3.3a的基函数绘制的二次曲线。
\end{minipage}
\end{center}

P3.3 端点插值: \(\mathbf{C}\left( 0\right)  = {\mathbf{P}}_{0}\) 和 \(\mathbf{C}\left( 1\right)  = {\mathbf{P}}_{n}\) ；

P3.4 仿射不变性\BookIndexMark{p0056}:通过对控制点应用仿射变换来对曲线应用仿射变换。令 \(\mathbf{r}\) 为 \({\mathcal{E}}^{3}\) 中的一个点(三维欧几里得空间)。仿射变换，记作 \(\Phi\) ，将 \({\mathcal{E}}^{3}\) 映射到 \({\mathcal{E}}^{3}\) ，其形式为

\[
\Phi \left( \mathbf{r}\right)  = A\mathbf{r} + \mathbf{v}
\]

其中 \(A\) 是一个 \(3 \times  3\) 矩阵，而 \(\mathbf{v}\) 是一个向量。仿射变换包括平移、旋转、缩放和剪切。B样条曲线的仿射不变性属性源于 \({N}_{i,p}\left( u\right)\) 的统一性划分属性。因此，设 \(\mathbf{r} = \sum {\alpha }_{i}{\mathbf{p}}_{i}\) ，其中 \({\mathbf{p}}_{i} \in  {\mathcal{E}}^{3}\) 和 \(\sum {\alpha }_{i} = 1\) 。然后

\[
\Phi \left( \mathbf{r}\right)  = \Phi \left( {\sum {\alpha }_{i}{\mathbf{p}}_{i}}\right)  = A\left( {\sum {\alpha }_{i}{\mathbf{p}}_{i}}\right)  + \mathbf{v} = \sum {\alpha }_{i}A{\mathbf{p}}_{i} + \sum {\alpha }_{i}\mathbf{v}
\]

\[
= \sum {\alpha }_{i}\left( {A{\mathbf{p}}_{i} + \mathbf{v}}\right)  = \sum {\alpha }_{i}\Phi \left( {\mathbf{p}}_{i}\right)
\]

P3.5 强凸包\BookIndexMark{p0057}属性:曲线包含在其控制多边形的凸包中；事实上，如果 \(u \in  \left\lbrack  {{u}_{i},{u}_{i + 1}}\right) ,p \leq  i < m - p - 1\) ，则 \(\mathbf{C}\left( u\right)\) 在控制点 \({\mathbf{P}}_{i - p},\ldots ,{\mathbf{P}}_{i}\) 的凸包中(图 3.4、3.5 和 3.6)。这源于 \({N}_{i,p}\left( u\right)\) 的非负性和统一性划分属性(属性 P2.3 和 P2.4)，以及当 \(u \in  \left\lbrack  {{u}_{i},{u}_{i + 1}}\right)\) 时 \({N}_{j,p}\left( u\right)  = 0\) 对于 \(j < i - p\) 和 \(j > i\) 的属性(属性 P2.2)。图 3.6 显示了如何构造一个包含直线段的二次曲线。由于 \({\mathbf{P}}_{2},{\mathbf{P}}_{3}\) 和 \({\mathbf{P}}_{4}\) 共线，强凸包属性强制曲线从 \(\mathbf{C}\left( {2/5}\right)\) 到 \(\mathbf{C}\left( {3/5}\right)\) 是一条直线段；

\begin{center}
\begin{minipage}{16cm}
\vspace*{.6\baselineskip}
\centering
\BookFigure[max width=\linewidth]{ch03/fig_3_04.tex}
\par\smallskip
图 3.4. 二次B样条曲线的强凸包性质；对于 \(u \in  \left\lbrack  {{u}_{i},{u}_{i + 1}}\right) ,\mathbf{C}\left( u\right)\) 位于三角形 \({\mathbf{P}}_{i - 2}{\mathbf{P}}_{i - 1}{\mathbf{P}}_{i}\) 内。
\end{minipage}
\end{center}

\begin{center}
\begin{minipage}{16cm}
\vspace*{.6\baselineskip}
\centering
\BookFigure[max width=\linewidth]{ch03/fig_3_05.tex}
\par\smallskip
图 3.5. 三次B样条曲线的强凸包性质；对于 \(u \in  \left\lbrack  {{u}_{i},{u}_{i + 1}}\right)\) ， \(\mathbf{C}\left( u\right)\) 位于四边形 \({\mathbf{P}}_{i - 3}{\mathbf{P}}_{i - 2}{\mathbf{P}}_{i - 1}{\mathbf{P}}_{i}\) 内。
\end{minipage}
\end{center}

\begin{center}
\begin{minipage}{16cm}
\vspace*{.6\baselineskip}
\centering
\BookFigure[max width=\linewidth]{ch03/fig_3_06.tex}
\par\smallskip
图 3.6. 在 \(U = \{ 0,0,0,1/5,2/5,3/5,4/5,1,1,1\}\) 上的二次B样条曲线。该曲线是 \(\mathbf{C}\left( {2/5}\right)\) 和 \(\mathbf{C}\left( {3/5}\right)\) 之间的直线。
\end{minipage}
\end{center}

P3.6 局部修改\BookIndexMark{p0058}方案:移动 \({\mathbf{P}}_{i}\) 仅在区间 \(\left\lbrack  {{u}_{i},{u}_{i + p + 1}}\right)\) 内改变 \(\mathbf{C}\left( u\right)\) (图 3.7)。这源于 \(u \notin  \left\lbrack  {{u}_{i},{u}_{i + p + 1}}\right)\) 的 \({N}_{i,p}\left( u\right)  = 0\) (属性 P2.1)。

\begin{center}
\begin{minipage}{16cm}
\vspace*{.6\baselineskip}
\centering
\BookFigure[max width=\linewidth]{ch03/fig_3_07.tex}
\par\smallskip
图 3.7. 在 \(U = \{ 0,0,0,0,1/4,1/2,3/4,1,1,1,1\}\) 上的三次曲线；移动 \({\mathbf{P}}_{4}\) (到 \({\mathbf{P}}_{4}^{\prime }\) )会改变区间 \(\left\lbrack  {1/4,1}\right)\) 中的曲线。
\end{minipage}
\end{center}

P3.7 控制多边形表示曲线的分段线性近似；通过节点插入或次数提升可以改善这种近似(见第5章)。一般来说，次数越低，B样条曲线越接近其控制多边形(见图3.8和3.9)。图3.9中的曲线使用相同的六个控制点定义，并且节点向量

\[
p = 1 : U = \left\{  {0,0,\frac{1}{5},\frac{2}{5},\frac{3}{5},\frac{4}{5},1,1}\right\}
\]

\[
p = 2 : U = \left\{  {0,0,0,\frac{1}{4},\frac{1}{2},\frac{3}{4},1,1,1}\right\}
\]

\[
p = 3 : U = \left\{  {0,0,0,0,\frac{1}{3},\frac{2}{3},1,1,1,1}\right\}
\]

\[
p = 4 : U = \left\{  {0,0,0,0,0,\frac{1}{2},1,1,1,1,1}\right\}
\]

\[
p = 5 : U = \{ 0,0,0,0,0,0,1,1,1,1,1,1\}
\]



这一现象的原因是直观的:次数越低，参与计算的控制点越少 \(\mathbf{C}\left( {u}_{0}\right)\) 对于任何给定的 \({u}_{0}\) 。极端情况是 \(p = 1\) ，在这种情况下，每个点 \(\mathbf{C}\left( u\right)\) 只是两个控制点之间的线性插值。在这种情况下，曲线就是控制多边形；

\begin{center}
\begin{minipage}{17cm}
\vspace*{.6\baselineskip}
\centering
\BookFigure[max width=\linewidth]{ch03/fig_3_08.tex}
\par\smallskip
图 3.8. B-spline 曲线。(a)在节点向量 \(U =\)  \(\{ 0,0,0,0,0,0,0,0,0,0,1,1,1,1,1,1,1,1,1,1\}\) 上的九次 Bézier 曲线；(b)使用在 \(U = \{ 0,0,0,1/8,2/8,3/8,4/8,5/8,6/8,7/8,1,1,1\}\) 上定义的相同控制多边形的二次曲线。
\end{minipage}
\end{center}

\begin{center}
\begin{minipage}{16cm}
\vspace*{.6\baselineskip}
\centering
\BookFigure[max width=\linewidth]{ch03/fig_3_09.tex}
\par\smallskip
图 3.9. 不同阶数的 B-spline 曲线，使用相同的控制多边形。
\end{minipage}
\end{center}

P3.8 沿着曲线从 \(u = 0\) 移动到 \(u = 1\) ， \({N}_{i,p}\left( u\right)\) 函数像开关一样起作用；当 \(u\) 经过一个节点时，一个 \({N}_{i,p}\left( u\right)\) (以及相应的 \({\mathbf{P}}_{i}\) )关闭，下一个则开启(图3.2和3.3)。





P3.9 变异减少性质\BookIndexMark{p0059}:没有平面与曲线的交点比与控制多边形的交点更多(对于二维曲线，将“平面”替换为“直线”) - 见 \cite{Lane83} 的证明；





P3.10 \(\mathbf{C}\left( u\right)\) 的连续性和可微性源于 \({N}_{i,p}\left( u\right)\) 的连续性和可微性(因为 \(\mathbf{C}\left( u\right)\) 只是 \({N}_{i,p}\left( u\right)\) 的线性组合)。因此， \(\mathbf{C}\left( u\right)\) 在结点区间的内部是无限可微的，并且在重数为 \(k\) 的结点处至少是 \(p - k\) 次连续可微的。图 3.10 显示了一条二次曲线 \(\left( {p = 2}\right)\) 。该曲线在重数为 1 的所有内部结点处是 \({\mathbf{C}}^{1}\) 连续的(第一导数是连续的，但第二导数不是)。在双重结点处， \(u = 4/5,\mathbf{C}\left( u\right)\) 仅是 \({\mathbf{C}}^{0}\) 连续的，因此存在一个尖点(视觉不连续性)。图 3.11 显示了一条在同一结点向量上定义的二次曲线。因此，这两条曲线使用相同的基函数 \({N}_{i,p}\left( u\right)\) 进行定义。但图 3.11 的曲线在 \(u = 4/5\) 处是 \({\mathbf{C}}^{1}\) 连续的；这并不明显，但可以通过第 3.3 节中给出的导数表达式来观察到。这只是因为不连续函数有时可以以某种方式组合，从而使结果是连续的。注意 \({\mathbf{P}}_{4}\) 、 \({\mathbf{P}}_{5}\) 和 \({\mathbf{P}}_{6}\) 是共线的，长度为 \(\left( {{\mathbf{P}}_{4}{\mathbf{P}}_{5}}\right)  =\) 长度 \(\left( {{\mathbf{P}}_{5}{\mathbf{P}}_{6}}\right)\) 。图 3.12 显示了一条在 \(u = 1/4\) 和 \(u = 1/2\) 处是 \({\mathbf{C}}^{2}\) 连续的三次曲线，但在双重结点 \(u = 3/4\) 处仅是 \({\mathbf{C}}^{1}\) 连续的。眼睛能够检测到第二导数的不连续性，但可能无法检测到第三及更高阶导数的不连续性。因此，三次曲线通常适用于视觉目的。



\begin{center}
\begin{minipage}{17cm}
\vspace*{.6\baselineskip}
\centering
\BookFigure[max width=\linewidth]{ch03/fig_3_10.tex}
\par\smallskip
图 3.10. 在 \(U = \{ 0,0,0,1/5,2/5,3/5,4/5,4/5,1,1,1\}\) 上的二次曲线，在 \(u = 4/5\) 处有一个尖点。
\end{minipage}
\end{center}





\begin{center}
\begin{minipage}{18cm}
\vspace*{.6\baselineskip}
\centering
\BookFigure[max width=\linewidth]{ch03/fig_3_11.tex}
\par\smallskip
图 3.11. 在 \(U = \{ 0,0,0,1/5,2/5,3/5,4/5,4/5,1,1,1\}\) 上的二次曲线；在 \(u = 4/5\) 处一阶导数是连续的。
\end{minipage}
\end{center}



\begin{center}
\begin{minipage}{18cm}
\vspace*{.6\baselineskip}
\centering
\BookFigure[max width=\linewidth]{ch03/fig_3_12.tex}
\par\smallskip
图 3.12. 在 \(U=\{0,0,0,0,1/4,1/2,3/4,3/4,1,1,1,1\}\) 上的三次曲线。\\ 在 \(u=1/4,1/2\) 处为 \(C^2\) 连续，在 \(u=3/4\) 处为 \(C^1\) 连续。
\end{minipage}
\end{center}

P3.11 使用多个(重合的\BookIndexMark{p0060})控制点是可能的(有时也是有用的)。图 3.13 显示了一个具有双重控制点的二次曲线， \({\mathbf{P}}_{2} = {\mathbf{P}}_{3}\) 。该曲线的有趣部分位于 \(\mathbf{C}\left( {1/4}\right)\) 和 \(\mathbf{C}\left( {3/4}\right)\) 之间。实际上， \(\mathbf{C}\left( {1/2}\right)  = {\mathbf{P}}_{2} = {\mathbf{P}}_{3}\) ，并且在 \(\mathbf{C}\left( {1/4}\right)\) 和 \(\mathbf{C}\left( {1/2}\right)\) 之间，以及 \(\mathbf{C}\left( {1/2}\right)\) 和 \(\mathbf{C}\left( {3/4}\right)\) 之间的曲线段都是直线。这是由于性质 P3.5 的结果，例如， \(\mathbf{C}\left( u\right)\) 在 \({\mathbf{P}}_{1}{\mathbf{P}}_{2}{\mathbf{P}}_{3}\) 的凸包(即一条线)内，如果 \(u \in  \left\lbrack  {1/4,1/2}\right)\) 。此外，由于节点 \(u = 1/2\) 的重数为 \(= 1\) ，因此曲线在该处必须是 \({\mathbf{C}}^{1}\) 连续的，尽管它有一个尖点(视觉不连续性)。这是因为在 \(u = 1/2\) 处一阶导数向量的大小连续趋向于零的结果。在下一节中，我们将看到在 \(u = 1/2\) 处的导数与差值 \({\mathbf{P}}_{3} - {\mathbf{P}}_{2}\) 成正比。图 3.14a 和 \({3.14}\mathrm{\;b}\) 是使用相同控制多边形的三次示例，包括双重控制点 \({\mathbf{P}}_{2} = {\mathbf{P}}_{3}\) ，但具有不同的节点向量。









\begin{center}
\begin{minipage}{18cm}
\vspace*{.6\baselineskip}
\centering
\BookFigure[max width=\linewidth]{ch03/fig_3_13.tex}
\par\smallskip
图 3.13. 在 \(U=\{0,0,0,1/4,1/2,3/4,1,1,1\}\) 上的二次曲线，具有重合控制点 \(\mathbf P_2=\mathbf P_3\)。
\end{minipage}
\end{center}

\begin{center}
\begin{minipage}{18cm}
\vspace*{.6\baselineskip}
\centering
\BookFigure[max width=\linewidth]{ch03/fig_3_14.tex}
\par\smallskip
图 3.14. 具有重合控制点 \(\mathbf P_2=\mathbf P_3\) 的三次曲线。\\ (a) \(U=\{0,0,0,0,1/4,3/4,1,1,1,1\}\)；(b) \(U=\{0,0,0,0,1/2,1/2,1,1,1,1\}\)。
\end{minipage}
\end{center}

\section{B样条曲线的导数}\BookIndexMark{p0061}

令 \({\mathbf{C}}^{\left( k\right) }\left( u\right)\) 表示 \(k\) 阶导数的 \(\mathbf{C}\left( u\right)\) 。如果 \(u\) 固定，我们可以通过计算基函数的 \(k\) 阶导数来获得 \({\mathbf{C}}^{\left( k\right) }\left( u\right)\) (见方程 [2.7]、[2.9] 和 [2.10] 以及算法 A2.3)。特别地

\[
{\mathbf{C}}^{\left( k\right) }\left( u\right)  = \mathop{\sum }\limits_{{i = 0}}^{n}{N}_{i,p}^{\left( k\right) }\left( u\right) {\mathbf{P}}_{i} \tag{3.3}
\]

考虑第 2.5 节的例子，设 \(p = 2,U = \{ 0,0,0,1,2,3,4,4,5,5,5\}\) 和 \(u = 5/2\) 。从方程 (2.7) 我们得到

\[
{N}_{2,2}^{\prime }\left( \frac{5}{2}\right)  = 0 - \frac{2}{3 - 1}\frac{1}{2} =  - \frac{1}{2}
\]

\[
{N}_{3,2}^{\prime }\left( \frac{5}{2}\right)  = \frac{2}{3 - 1}\frac{1}{2} - \frac{2}{4 - 2}\frac{1}{2} = 0
\]

\[
{N}_{4,2}^{\prime }\left( \frac{5}{2}\right)  = \frac{2}{4 - 2}\frac{1}{2} - 0 = \frac{1}{2}
\]

由此可得

\[
{\mathbf{C}}^{\prime }\left( \frac{5}{2}\right)  =  - \frac{1}{2}{\mathbf{P}}_{2} + \frac{1}{2}{\mathbf{P}}_{4}
\]





一个算法用于计算 B样条曲线上的点及其所有导数，直到并包括在固定的 \(u\) 值下的 \(d\) 阶导数。我们允许 \(d > p\) ，尽管在这种情况下导数为 0(对于非有理曲线)；这些导数对于有理曲线是必要的。算法的输入为 \(u,d\) ，以及 B样条曲线，通过本书其余部分定义为

n : 控制点的数量为 \(n + 1\) ;

p : 曲线的阶数；

U : 节点；

P : 控制点。

输出为数组 \(\mathrm{{CK}}\left\lbrack  \right\rbrack\) ，其中 \(\mathrm{{CK}}\left\lbrack  \mathrm{k}\right\rbrack\) 是 \(k\) 阶导数 \(0 \leq  k \leq  d\) 。我们使用算法 A2.1 和 A2.3。一个局部数组 nders[][] 用于存储基函数的导数。

% book-code: algorithm A3.2
\par\medskip\noindent\textbf{算法 A3.2（伪代码）}\BookIndexMark{p0062}\par
\begin{minted}[escapeinside=@@,mathescape]{c}
CurveDerivsAlg1(n, p, U, P, u, d, CK)
{ /* Compute curve derivatives */
    /* Input: n, p, U, P, u, d */
    /* Output: CK */
    du = min(d, p);
    for (k=p+1; k<=d; k++) CK[k] = 0.0;
    span = FindSpan(n, p, u, U);
    DersBasisFuns(span, u, p, du, U, nders);
    for ( k = 0 ; k <= du ; k ++ )
    {
        CK[k] = 0.0 ;
        for (j=0; j<=p; j++)
            CK[k] = CK[k] + nders[k][j]*P[span-p+j];
    }
}
\end{minted}

现在，我们想要正式区分 \(p\) 阶 B-样条曲线，而不是修正 \(u\) 。

样条曲线，

\[
\mathbf{C}\left( u\right)  = \mathop{\sum }\limits_{{i = 0}}^{n}{N}_{i,p}\left( u\right) {\mathbf{P}}_{i}
\]

定义在节点向量上

\[
U = \{ \underset{p + 1}{\underbrace{0,\ldots ,0}},{u}_{p + 1},\ldots ,{u}_{m - p - 1},\underset{p + 1}{\underbrace{1,\ldots ,1}}\}
\]

从方程 (3.3) 和 (2.7) 我们得到

\[
{\mathbf{C}}^{\prime }\left( u\right)  = \mathop{\sum }\limits_{{i = 0}}^{n}{N}_{i,p}^{\prime }\left( u\right) {\mathbf{P}}_{i}
\]

\[
= \mathop{\sum }\limits_{{i = 0}}^{n}\left( {\frac{p}{{u}_{i + p} - {u}_{i}}{N}_{i,p - 1}\left( u\right)  - \frac{p}{{u}_{i + p + 1} - {u}_{i + 1}}{N}_{i + 1,p - 1}\left( u\right) }\right) {\mathbf{P}}_{i}
\]

\[
= \left( {p\mathop{\sum }\limits_{{i =  - 1}}^{{n - 1}}{N}_{i + 1,p - 1}\left( u\right) \frac{{\mathbf{P}}_{i + 1}}{{u}_{i + p + 1} - {u}_{i + 1}}}\right)
\]

94 B-样条曲线和曲面

\[
- \left( {p\mathop{\sum }\limits_{{i = 0}}^{n}{N}_{i + 1,p - 1}\left( u\right) \frac{{\mathbf{P}}_{i}}{{u}_{i + p + 1} - {u}_{i + 1}}}\right)
\]

\[
= p\frac{{N}_{0,p - 1}\left( u\right) {\mathbf{P}}_{0}}{{u}_{p} - {u}_{0}} + p\mathop{\sum }\limits_{{i = 0}}^{{n - 1}}{N}_{i + 1,p - 1}\left( u\right) \frac{\left( {\mathbf{P}}_{i + 1} - {\mathbf{P}}_{i}\right) }{{u}_{i + p + 1} - {u}_{i + 1}} - p\frac{{N}_{n + 1,p - 1}\left( u\right) {\mathbf{P}}_{n}}{{u}_{n + p + 1} - {u}_{n + 1}}
\]

第一个和最后一个项的值为 \(0/0\) ，根据定义为 0。因此

\[
{\mathbf{C}}^{\prime }\left( u\right)  = p\mathop{\sum }\limits_{{i = 0}}^{{n - 1}}{N}_{i + 1,p - 1}\left( u\right) \frac{\left( {\mathbf{P}}_{i + 1} - {\mathbf{P}}_{i}\right) }{{u}_{i + p + 1} - {u}_{i + 1}} = \mathop{\sum }\limits_{{i = 0}}^{{n - 1}}{N}_{i + 1,p - 1}\left( u\right) {\mathbf{Q}}_{i}
\]

其中

\[
{\mathbf{Q}}_{i} = p\frac{{\mathbf{P}}_{i + 1} - {\mathbf{P}}_{i}}{{u}_{i + p + 1} - {u}_{i + 1}} \tag{3.4}
\]

现在设 \({U}^{\prime }\) 为通过从 \(U\) 中去掉第一个和最后一个节点得到的节点向量，即

\[
{U}^{\prime } = \{ \underset{p}{\underbrace{0,\ldots ,0}},{u}_{p + 1},\ldots ,{u}_{m - p - 1},\underset{p}{\underbrace{1,\ldots ,1}}\}  \tag{3.5}
\]

\(\left( {U}^{\prime }\right.\) 有 \(m - 1\) 个节点)。然后很容易检查在 \(U\) 上计算的函数 \({N}_{i + 1,p - 1}\left( u\right)\) 等于在 \({U}^{\prime }\) 上计算的 \({N}_{i,p - 1}\left( u\right)\) 。因此

\[
{\mathbf{C}}^{\prime }\left( u\right)  = \mathop{\sum }\limits_{{i = 0}}^{{n - 1}}{N}_{i,p - 1}\left( u\right) {\mathbf{Q}}_{i} \tag{3.6}
\]

其中 \({\mathbf{Q}}_{i}\) 由方程 (3.4) 定义，而 \({N}_{i,p - 1}\left( u\right)\) 在 \({U}^{\prime }\) 上计算。因此， \({\mathbf{C}}^{\prime }\left( u\right)\) 是一个 (p - 1) 阶 B-样条曲线。

\section*{示例}

Ex3.1 设 \(\mathbf{C}\left( u\right)  = \mathop{\sum }\limits_{{i = 0}}^{4}{N}_{i,2}\left( u\right) {\mathbf{P}}_{i}\) 为定义在

\[
U = \{ 0,0,0,2/5,3/5,1,1,1\}
\]

然后 \({U}^{\prime } = \{ 0,0,2/5,3/5,1,1\}\) 和 \({\mathbf{C}}^{\prime }\left( u\right)  = \mathop{\sum }\limits_{{i = 0}}^{3}{N}_{i,1}\left( u\right) {\mathbf{Q}}_{i}\) ，其中

\[
\mathbf Q_{0}=\frac{2\left(\mathbf P_{1}-\mathbf P_{0}\right)}{\frac25-0}=5\left(\mathbf P_{1}-\mathbf P_{0}\right)
\]

\[
\mathbf Q_{1}=\frac{2\left(\mathbf P_{2}-\mathbf P_{1}\right)}{\frac35-0}=\frac{10}{3}\left(\mathbf P_{2}-\mathbf P_{1}\right)
\]

\[
\mathbf Q_{2}=\frac{2\left(\mathbf P_{3}-\mathbf P_{2}\right)}{1-\frac25}=\frac{10}{3}\left(\mathbf P_{3}-\mathbf P_{2}\right)
\]

\[
\mathbf Q_{3}=\frac{2\left(\mathbf P_{4}-\mathbf P_{3}\right)}{1-\frac35}=5\left(\mathbf P_{4}-\mathbf P_{3}\right)
\]

\(\mathbf{C}\left( u\right)\) 和 \({\mathbf{C}}^{\prime }\left( u\right)\) 分别如图 \({3.15}\mathrm{a}\) 和 \({3.15}\mathrm{\;b}\) 所示。

\begin{center}
\begin{minipage}{18cm}
\vspace*{.6\baselineskip}
\centering
\BookFigure[max width=\linewidth]{ch03/fig_3_15.tex}
\par\smallskip
图 3.15. (a) \(U=\{0,0,0,2/5,3/5,1,1,1\}\) 上的二次曲线；切向量按 \(1/2\) 缩放显示。\\ (b) \(U'=\{0,0,2/5,3/5,1,1\}\) 上的一次导数曲线，控制点由式(3.4)计算。
\end{minipage}
\end{center}


Ex3.2 设 \(\mathbf{C}\left( u\right)  = \mathop{\sum }\limits_{{i = 0}}^{6}{N}_{i,3}\left( u\right) {\mathbf{P}}_{i}\) 为定义在

\[
U = \{ 0,0,0,0,2/5,3/5,3/5,1,1,1,1\}
\]

然后 \({U}^{\prime } = \left\{  {0,0,0,2/5,3/5,3/5,1,1,1}\right\}\) 和 \({\mathbf{C}}^{\prime }\left( u\right)  = \mathop{\sum }\limits_{{i = 0}}^{5}{N}_{i,2}\left( u\right) {\mathbf{Q}}_{i}\) ，

其中

\[
\mathbf Q_{0}=\frac{3\left(\mathbf P_{1}-\mathbf P_{0}\right)}{\frac25-0}=\frac{15}{2}\left(\mathbf P_{1}-\mathbf P_{0}\right)
\]

\[
\mathbf Q_{1}=\frac{3\left(\mathbf P_{2}-\mathbf P_{1}\right)}{\frac35-0}=5\left(\mathbf P_{2}-\mathbf P_{1}\right)
\]

\[
\mathbf Q_{2}=\frac{3\left(\mathbf P_{3}-\mathbf P_{2}\right)}{\frac35-0}=5\left(\mathbf P_{3}-\mathbf P_{2}\right)
\]

\[
\mathbf Q_{3}=\frac{3\left(\mathbf P_{4}-\mathbf P_{3}\right)}{1-\frac25}=5\left(\mathbf P_{4}-\mathbf P_{3}\right)
\]

\[
\mathbf Q_{4}=\frac{3\left(\mathbf P_{5}-\mathbf P_{4}\right)}{1-\frac35}=\frac{15}{2}\left(\mathbf P_{5}-\mathbf P_{4}\right)
\]

\[
\mathbf Q_{5}=\frac{3\left(\mathbf P_{6}-\mathbf P_{5}\right)}{1-\frac35}=\frac{15}{2}\left(\mathbf P_{6}-\mathbf P_{5}\right)
\]

\(\mathbf{C}\left( u\right)\) 和 \({\mathbf{C}}^{\prime }\left( u\right)\) 分别如图 3.16a 和 \({3.16}\mathrm{\;b}\) 所示。注意 \({\mathbf{C}}^{\prime }\left( u\right)\) 是一个在双结点 \(u = 3/5\) 处具有尖点的二次曲线。

\begin{center}
\begin{minipage}{18cm}
\vspace*{.6\baselineskip}
\centering
\BookFigure[max width=\linewidth]{ch03/fig_3_16.tex}
\par\medskip
\BookFigure[max width=\linewidth]{ch03/fig_3_17.tex}
\par\smallskip
图 3.16. (a) \(U=\{0,0,0,0,2/5,3/5,3/5,1,1,1,1\}\) 上的三次曲线及切向量。\\ (b) \(U'=\{0,0,0,2/5,3/5,3/5,1,1,1\}\) 上的二次曲线 \(\mathbf C'(u)\)；两图的导数量按 \(1/3\) 缩放显示。
\end{minipage}
\end{center}


Ex3.3 回想一下， \(p\) 阶 Bézier 曲线是 \(U = \{ 0,\ldots ,0,1,\ldots ,1\}\) 上的 B-spline 曲线(没有内部结点)，方程 (3.4) 简化为 \({\mathbf{Q}}_{i} =\)  \(p\left( {{\mathbf{P}}_{i + 1} - {\mathbf{P}}_{i}}\right)\) 对于 \(0 \leq  i \leq  n - 1\) 。由于 \(n = p\) 和 \({N}_{i,p - 1}\left( u\right)  = {B}_{i,n - 1}\left( u\right)\) ，伯恩斯坦多项式，方程 (3.6) 等价于方程 (1.9)。







B样条曲线在端点的第一导数由以下公式给出

\[
{\mathbf{C}}^{\prime }\left( 0\right)  = {\mathbf{Q}}_{0} = \frac{p}{{u}_{p + 1}}\left( {{\mathbf{P}}_{1} - {\mathbf{P}}_{0}}\right)
\]

\[
{\mathbf{C}}^{\prime }\left( 1\right)  = {\mathbf{Q}}_{n - 1} = \frac{p}{1 - {u}_{m - p - 1}}\left( {{\mathbf{P}}_{n} - {\mathbf{P}}_{n - 1}}\right)  \tag{3.7}
\]

(见示例Ex3.1和Ex3.2，以及图3.15(a)和(b)以及图3.16(a)和(b))。注意，在图3.15b和 \({3.16}\mathrm{\;b}\) 中，导数向量和控制点差异被缩小以便于更好的可视化，分别由 \(1/2\) 和 \(1/3\) 表示。

由于 \({\mathbf{C}}^{\prime }\left( u\right)\) 是B样条曲线，我们递归地应用方程(3.4)至(3.6)以获得更高阶的导数。设 \({\mathbf{P}}_{i}^{\left( 0\right) } = {\mathbf{P}}_{i}\) ，我们写成

\[
\mathbf{C}\left( u\right)  = {\mathbf{C}}^{\left( 0\right) }\left( u\right)  = \mathop{\sum }\limits_{{i = 0}}^{n}{N}_{i,p}\left( u\right) {\mathbf{P}}_{i}^{\left( 0\right) }
\]

然后

\[
{\mathbf{C}}^{\left( k\right) }\left( u\right)  = \mathop{\sum }\limits_{{i = 0}}^{{n - k}}{N}_{i,p - k}\left( u\right) {\mathbf{P}}_{i}^{\left( k\right) } \tag{3.8}
\]

具有

\[
{\mathbf{P}}_{i}^{\left( k\right) } = \left\{  \begin{array}{ll} {\mathbf{P}}_{i} & k = 0 \\  \frac{p - k + 1}{{u}_{i + p + 1} - {u}_{i + k}}\left( {{\mathbf{P}}_{i + 1}^{\left( k - 1\right) } - {\mathbf{P}}_{i}^{\left( k - 1\right) }}\right) & k > 0 \end{array}\right.
\]



和

\[
{U}^{\left( k\right) } = \left\{  {\underset{p - k + 1}{\underbrace{0,\ldots ,0}},{u}_{p + 1},\ldots ,{u}_{m - p - 1},\underset{p - k + 1}{\underbrace{1,\ldots ,1}}}\right\}
\]

算法A3.3是方程(3.8)的非递归实现。它计算所有导数曲线的控制点，直到并包括 \(d\) 阶导数 \(\left( {d \leq  p}\right)\) 。输出时，PK[k][i]是 \(k\) 阶导数曲线的 \(i\) 阶控制点，其中 \(0 \leq  k \leq  d\) 和 \({r}_{1} \leq  i \leq  {r}_{2} - k\) 。如果 \({r}_{1} = 0\) 和 \({r}_{2} = n\) ，则计算所有控制点。

% book-code: algorithm A3.3
\par\medskip\noindent\textbf{算法 A3.3（伪代码）}\BookIndexMark{p0063}\par
\begin{minted}[escapeinside=@@,mathescape]{c}
CurveDerivCpts(n, p, U, P, d, r1, r2, PK)
{ /* Compute control points of curve derivatives */
    /* Input: n, p, U, P, d, r1, r2 */
    /* Output: PK */
    r = r2-r1;
    for (i = 0; i<=r; i++)
        PK[0][i] = P[r1+i];
    for (k=1; k<=d; k++)
    {
        tmp = p-k+1;
        for (i = 0; i<=r-k; i++)
            PK[k][i] = tmp*(PK[k-1][i+1]-PK[k-1][i])/
        ( U[r1 + i + p + 1]- U[r1 + i + k]) ;
    }
}
\end{minted}

使用方程(3.8)，我们计算B样条曲线 \(\left( {p > 1}\right)\) 在端点 \(u = 0\) 处的第二导数

\[
{\mathbf{C}}^{\left( 2\right) }\left( 0\right)  = {\mathbf{P}}_{0}^{\left( 2\right) } = \frac{p - 2 + 1}{{u}_{p + 1} - {u}_{2}}\left( {{\mathbf{P}}_{1}^{\left( 1\right) } - {\mathbf{P}}_{0}^{\left( 1\right) }}\right)
\]

\[
= \frac{p - 1}{{u}_{p + 1} - {u}_{2}}\left\lbrack  {\frac{p}{{u}_{p + 2} - {u}_{2}}\left( {{\mathbf{P}}_{2}^{\left( 0\right) } - {\mathbf{P}}_{1}^{\left( 0\right) }}\right)  - \frac{p}{{u}_{p + 1} - {u}_{1}}\left( {{\mathbf{P}}_{1}^{\left( 0\right) } - {\mathbf{P}}_{0}^{\left( 0\right) }}\right) }\right\rbrack
\]

从 \({u}_{1} = {u}_{2} = 0\) 可以得出

\[
{\mathbf{C}}^{\left( 2\right) }\left( 0\right)  = \frac{p\left( {p - 1}\right) }{{u}_{p + 1}}\left\lbrack  {\frac{{\mathbf{P}}_{0}}{{u}_{p + 1}} - \frac{\left( {{u}_{p + 1} + {u}_{p + 2}}\right) {\mathbf{P}}_{1}}{{u}_{p + 1}{u}_{p + 2}} + \frac{{\mathbf{P}}_{2}}{{u}_{p + 2}}}\right\rbrack   \tag{3.9}
\]

类似地，

\[
{\mathbf{C}}^{\left( 2\right) }\left( 1\right)  = \frac{p\left( {p - 1}\right) }{1 - {u}_{m - p - 1}} \times
\]

\[
\left\lbrack  {\frac{{\mathbf{P}}_{n}}{1 - {u}_{m - p - 1}} - \frac{\left( {2 - {u}_{m - p - 1} - {u}_{m - p - 2}}\right) {\mathbf{P}}_{n - 1}}{\left( {1 - {u}_{m - p - 1}}\right) \left( {1 - {u}_{m - p - 2}}\right) } + \frac{{\mathbf{P}}_{n - 2}}{1 - {u}_{m - p - 2}}}\right\rbrack   \tag{3.10}
\]

注意，对于贝塞尔曲线，这些方程简化为方程 (1.10) 的相应表达式。图 3.17 显示了图 3.15a 的二次曲线及其向量 \({\mathbf{C}}^{\left( 2\right) }\left( 0\right)\) 和 \({\mathbf{C}}^{\left( 2\right) }\left( 1\right)\) 。 \({\mathbf{C}}^{\left( 2\right) }\left( u\right)\) 是一条分段的零度曲线，即在三个区间 \(\lbrack 0,2/5)\) 、 \(\left\lbrack  {2/5,3/5}\right)\) 和 \(\left\lbrack  {3/5,1}\right\rbrack\) 上，它是一个常量(但不同)的向量。

\begin{center}
\begin{minipage}{16cm}
\vspace*{.6\baselineskip}
\centering
\BookFigure[max width=\linewidth]{ch03/fig_3_18.tex}
\par\smallskip
图 3.17. 图3.15a曲线端点的二阶导数向量，按 \(1/10\) 缩放显示。
\end{minipage}
\end{center}

我们用另一个算法结束本节，该算法用于计算 B 样条曲线上的一点及其所有导数，直到并包括在固定的 \(u\) 值下的 \(d\) 阶导数(与算法 A3.2 比较)。该算法基于方程 (3.8) 和算法 A3.3。我们假设有一个例程 AllBasisFuns，这是对 BasisFuns(算法 A2.2)的简单修改，用于返回从 0 到 \(p\) 的所有非零基函数。特别地， \(\mathrm{N}\left\lbrack  j\right\rbrack  \left\lbrack  i\right\rbrack\) 是 \(i\) 阶基函数 \({N}_{\text{ span } - i + j,i}\left( u\right)\) 的值，其中 \(0 \leq  i \leq  p\) 和 \(0 \leq  j \leq  i\) 。

% book-code: algorithm A3.4
\par\medskip\noindent\textbf{算法 A3.4（伪代码）}\BookIndexMark{p0064}\par
\begin{minted}[escapeinside=@@,mathescape]{c}
CurveDerivsAlg2(n, p, U, P, u, d, CK)
{ /* Compute curve derivatives */
    /* Input: n, p, U, P, u, d */
    /* Output: CK */
    du = min(d, p);
    for ( k = p + 1;k <= d;k ++ ) CK[k]= 0.0 ;
    span = FindSpan( n,p,u,U) ;
    AllBasisFuns(span, u, p, U, N);
    CurveDerivCpts(n, p, U, P, du, span-p, span, PK);
    for ( k = 0;k <= du;k ++ )
    {
        CK[k]= 0.0 ;
        for ( j = 0;j <= p - k;j ++ )
            CK[k]= CK[k]+ N[j][p - k]* PK[k][j];
    }
}
\end{minted}

图 3.18 显示了一条在 \(u = 2/5\) 处计算的一阶、二阶和三阶导数的三次曲线。(导数被缩放至 \(2/5\) 。)



\begin{center}
\begin{minipage}{18cm}
\vspace*{.6\baselineskip}
\centering
\BookFigure[max width=\linewidth]{ch03/fig_3_19.tex}
\par\smallskip
图 3.18. 在 \(U=\{0,0,0,0,1/4,3/4,1,1,1,1\}\) 上的三次曲线。\\ 在 \(u=2/5\) 处计算的一至三阶导数向量，均按 \(2/5\) 缩放显示。
\end{minipage}
\end{center}

\section{B样条曲面的定义与性质}\BookIndexMark{p0065}

B 样条曲面是通过取一个双向控制点网、两个节点向量以及单变量 B 样条函数的乘积获得的。

\[
\mathbf{S}\left( {u,v}\right)  = \mathop{\sum }\limits_{{i = 0}}^{n}\mathop{\sum }\limits_{{j = 0}}^{m}{N}_{i,p}\left( u\right) {N}_{j,q}\left( v\right) {\mathbf{P}}_{i,j} \tag{3.11}
\]



具有

\[
U = \{ \underset{p + 1}{\underbrace{0,\ldots ,0}},{u}_{p + 1},\ldots ,{u}_{r - p - 1},\underset{p + 1}{\underbrace{1,\ldots ,1}}\}
\]

\[
V = \left\{  {\underset{q + 1}{\underbrace{0,\ldots ,0}},{v}_{q + 1},\ldots ,{v}_{s - q - 1},\underset{q + 1}{\underbrace{1,\ldots ,1}}}\right\}
\]

\(U\) 具有 \(r + 1\) 个节点，且 \(V\) 具有 \(s + 1\) 个节点。方程 (3.2) 的形式为

\[
r = n + p + 1\text{ and }s = m + q + 1 \tag{3.12}
\]

设 \(U\) 及三次基函数 \(N_{i,3}(u)\) 如图3.2a所示，二次基函数 \(N_{j,2}(v)\) 定义在 \(V=\{0,0,0,1/5,2/5,3/5,3/5,4/5,1,1,1\}\) 上。图3.19a、b分别显示张量积基函数 \(N_{4,3}(u)N_{4,2}(v)\) 和 \(N_{4,3}(u)N_{2,2}(v)\)。图3.20～3.25显示B样条曲面的示例。

\begin{center}
\begin{minipage}{18cm}
\vspace*{.6\baselineskip}
\centering
\BookFigure[max width=\linewidth]{ch03/fig_3_20.tex}
\par\smallskip
图 3.19. 三次 \(\times\) 二次张量积基函数：(a) \(N_{4,3}(u)N_{4,2}(v)\)；(b) \(N_{4,3}(u)N_{2,2}(v)\)。\\ \(U=\{0,0,0,0,1/4,1/2,3/4,1,1,1,1\}\)，\\ \(V=\{0,0,0,1/5,2/5,3/5,3/5,4/5,1,1,1\}\)。
\end{minipage}
\end{center}

计算固定 (u, v) 参数值下 B-spline 曲面上的一个点需要五个步骤:

1. 找到 \(u\) 所在的节点区间，记为 \(u \in  \left\lbrack  {{u}_{i},{u}_{i + 1}}\right)\) (算法 A2.1)；

2. 计算非零基函数 \({N}_{i - p,p}\left( u\right) ,\ldots ,{N}_{i,p}\left( u\right)\) (A2.2)；

3. 找到 \(v\) 所在的节点区间，记为 \(v \in  \left\lbrack  {{v}_{j},{v}_{j + 1}}\right)\) (A2.1)；

4. 计算非零基函数 \({N}_{j - q,q}\left( v\right) ,\ldots ,{N}_{j,q}\left( v\right)\) (A2.2)；



5. 将非零基函数的值与相应的控制点相乘。

最后一步的形式为

\[
\mathbf{S}\left( {u,v}\right)  = {\left\lbrack  {N}_{k,p}\left( u\right) \right\rbrack  }^{T}\left\lbrack  {\mathbf{P}}_{k,l}\right\rbrack  \left\lbrack  {{N}_{l,q}\left( v\right) }\right\rbrack  \;i - p \leq  k \leq  i,\;j - q \leq  l \leq  j \tag{3.13}
\]

注意 \({\left\lbrack  {N}_{k,p}\left( u\right) \right\rbrack  }^{T}\) 是一个 1 \(\times  \left( {p + 1}\right)\) 行的标量向量， \(\left\lbrack  {\mathbf{P}}_{k,l}\right\rbrack\) 是一个 \(\left( {p + 1}\right)  \times  \left( {q + 1}\right)\) 的控制点矩阵，而 \(\left\lbrack  {{N}_{l,q}\left( v\right) }\right\rbrack\) 是一个 \(\left( {q + 1}\right)  \times  1\) 列的标量向量。

\section*{示例}

Ex3.4 设 \(p = q = 2\) 和 \(\mathop{\sum }\limits_{{i = 0}}^{4}\mathop{\sum }\limits_{{j = 0}}^{5}{N}_{i,2}\left( u\right) {N}_{j,2}\left( v\right) {\mathbf{P}}_{i,j}\) ，并且

\[
U = \{ 0,0,0,2/5,3/5,1,1,1\}
\]

\[
V = \{ 0,0,0,1/5,1/2,4/5,1,1,1\}
\]

计算 \(\mathbf{S}\left( {1/5,3/5}\right)\) 。然后 \(1/5 \in  \left\lbrack  {{u}_{2},{u}_{3}}\right)\) 和 \(3/5 \in  \left\lbrack  {{v}_{4},{v}_{5}}\right)\) ，并且

\[
\mathbf{S}\left( {\frac{1}{5},\frac{3}{5}}\right)  = \left\lbrack  \begin{array}{lll} {N}_{0,2}\left( \frac{1}{5}\right) & {N}_{1,2}\left( \frac{1}{5}\right) & {N}_{2,2}\left( \frac{1}{5}\right)  \end{array}\right\rbrack
\]

\[
\left\lbrack  \begin{array}{lll} {\mathbf{P}}_{0,2} & {\mathbf{P}}_{0,3} & {\mathbf{P}}_{0,4} \\  {\mathbf{P}}_{1,2} & {\mathbf{P}}_{1,3} & {\mathbf{P}}_{1,4} \\  {\mathbf{P}}_{2,2} & {\mathbf{P}}_{2,3} & {\mathbf{P}}_{2,4} \end{array}\right\rbrack  \left\lbrack  \begin{array}{l} {N}_{2,2}\left( \frac{3}{5}\right) \\  {N}_{3,2}\left( \frac{3}{5}\right) \\  {N}_{4,2}\left( \frac{3}{5}\right)  \end{array}\right\rbrack
\]

算法 A3.5 计算在固定 (u, v) 值下的 B 样条曲面上的点。为了提高效率，它使用一个局部数组 temp[] 来存储向量/矩阵乘积 \({\left\lbrack  {N}_{k,p}\left( u\right) \right\rbrack  }^{T}\left\lbrack  {\mathbf{P}}_{k,l}\right\rbrack\) 。然后，将生成的点向量(在 temp[] 中)进行乘法运算

与向量 \(\left\lbrack  {{N}_{l,q}\left( v\right) }\right\rbrack\) 。

% book-code: algorithm A3.5
\par\medskip\noindent\textbf{算法 A3.5（伪代码）}\BookIndexMark{p0066}\par
\begin{minted}[escapeinside=@@,mathescape]{c}
SurfacePoint( n,p,U,m,q,V,P,u,v,S)
{ /* Compute surface point */
    /* Input: n, p, U, m, q, V, P, u, v */
    /* Output: S */
    uspan = FindSpan(n, p, u, U);
    BasisFuns(uspan, u, p, U, Nu);
    vspan = FindSpan(m, q, v, V);
    BasisFuns(vspan, v, q, V, Nv);
    for (l=0; l<=q; l++)
    {
        temp[1] = 0.0 ;
        for ( k = 0;k <= p;k ++ )
            temp[1] = temp[1] + Nu[k]*P[uspan-p+k][vspan-q+1];
    }
    s = 0.0 ;
    for (l=0; l<=q; l++)
        S = S + Nv[l]*temp[l];
}
\end{minted}



\begin{center}
\begin{minipage}{18cm}
\vspace*{.6\baselineskip}
\centering
\BookFigure[max width=\linewidth]{ch03/fig_3_21.tex}
\par\medskip
\BookFigure[max width=\linewidth]{ch03/fig_3_22.tex}
\par\smallskip
图 3.20. 按代表性控制点构造的B样条曲面：(a) \(5\times4\) 控制网；(b) 对应的双三次曲面。\\ \(p=q=3,\quad U=\{0,0,0,0,1/2,1,1,1,1\}\)，\\ \(V=\{0,0,0,0,1,1,1,1\}\)。
\end{minipage}
\end{center}

张量积基函数的性质源于第 2 章中列出的单变量基函数的相应性质。

P3.12 非负性: \({N}_{i,p}\left( u\right) {N}_{j,q}\left( v\right)  \geq  0\) 对所有 \(i,j,p,q,u,v\) ；

P3.13\BookIndexMark{p0067} 单位划分: \(\mathop{\sum }\limits_{{i = 0}}^{n}\mathop{\sum }\limits_{{j = 0}}^{m}{N}_{i,p}\left( u\right) {N}_{j,q}\left( v\right)  = 1\) 对所有 \(\left( {u,v}\right)  \in  \left\lbrack  {0,1}\right\rbrack   \times\)  \(\left\lbrack  {0,1}\right\rbrack\) ；

P3.14 如果 \(n = p,m = q,U = \{ 0,\ldots ,0,1,\ldots ,1\}\) ，并且 \(V = \{ 0,\ldots ,0,1,\ldots ,1\}\) ，则对于所有 \(i,j\) ，有 \({N}_{i,p}\left( u\right) {N}_{j,q}\left( v\right)  = {B}_{i,n}\left( u\right) {B}_{j,m}\left( v\right)\) ；即，B 样条函数的乘积退化为伯恩斯坦多项式的乘积；

P3.15\BookIndexMark{p0068} \({N}_{i,p}\left( u\right) {N}_{j,q}\left( v\right)  = 0\) 如果 (u, v) 在矩形 \(\left\lbrack  {{u}_{i},{u}_{i + p + 1}}\right)  \times\)  \(\left\lbrack  {{v}_{j},{v}_{j + q + 1}}\right)\) 之外(见图 \({3.19}\mathrm{a}\) 和 \({3.19}\mathrm{\;b}\) )；

P3.16 在任何给定的矩形中， \(\left\lbrack  {{u}_{{i}_{0}},{u}_{{i}_{0} + 1}}\right)  \times  \left\lbrack  {{v}_{{j}_{0}},{v}_{{j}_{0} + 1}}\right)\) ，至多有 \(\left( {p + 1}\right) \left( {q + 1}\right)\) 个基函数非零，特别是 \({N}_{i,p}\left( u\right) {N}_{j,q}\left( v\right)\) 对于 \({i}_{0} - p \leq\)  \(i \leq  {i}_{0}\) 和 \({j}_{0} - q \leq  j \leq  {j}_{0}\) ；

P3.17 如果 \(p > 0\) 和 \(q > 0\) ，那么 \({N}_{i,p}\left( u\right) {N}_{j,q}\left( v\right)\) 恰好达到一个最大值(见图 3.19a 和 3.19b)；

在由 \(u\) 和 \(v\) 节点线形成的矩形内部，函数是一个双变量多项式，所有 \({N}_{i,p}\left( u\right)\)  \({N}_{j,q}\left( v\right)\) 的偏导数存在；在 \(u\) 节点( \(v\) 节点)处，它在 \(u\left( v\right)\) 方向上是 \(p - k\left( {q - k}\right)\) 次可微的，其中 \(k\) 是节点的重数。在图 3.19a 中， \({N}_{4,3}\left( u\right) {N}_{4,2}\left( v\right)\) 关于 \(v\) 的第一个偏导数在节点线 \(v = 3/5\) 上是不连续的，那里 \({N}_{4,2}\left( v\right)\) 有一个尖点。关于 \(u\) 的第二个偏导数处处连续，因为 \({N}_{4,3}\left( u\right)\) 是 \({\mathbf{C}}^{2}\) 连续的。

B-spline 曲面具有以下特性:

P3.19 如果 \(n = p,m = q,U = \{ 0,\ldots ,0,1,\ldots ,1\}\) ，并且 \(V = \{ 0,\ldots ,0,1,\ldots ,1\}\) ，那么 \(\mathbf{S}\left( {u,v}\right)\) 是一个 Bézier 曲面；这来自 P3.14；

P3.20 该曲面插值四个角控制点: \(\mathbf{S}\left( {0,0}\right)  = {\mathbf{P}}_{0,0}\) ， \(\mathbf{S}\left( {1,0}\right)  = {\mathbf{P}}_{n,0},\mathbf{S}\left( {0,1}\right)  = {\mathbf{P}}_{0,m}\) 和 \(\mathbf{S}\left( {1,1}\right)  = {\mathbf{P}}_{n,m}\) (见图 3.20 至 3.25)；这来自 P3.13 和该恒等式

\[
{N}_{0,p}\left( 0\right) {N}_{0,q}\left( 0\right)  = {N}_{n,p}\left( 1\right) {N}_{0,q}\left( 0\right)  = {N}_{0,p}\left( 0\right) {N}_{m,q}\left( 1\right)
\]

\[
= {N}_{n,p}\left( 1\right) {N}_{m,q}\left( 1\right)  = 1
\]

P3.21 仿射不变性:通过对控制点应用仿射变换来对曲面应用该变换；这来自 P3.13；

P3.22 强凸包性质:如果 \(\left( {u,v}\right)  \in  \left\lbrack  {{u}_{{i}_{0}},{u}_{{i}_{0} + 1}}\right)  \times  \left\lbrack  {{v}_{{j}_{0}},{v}_{{j}_{0} + 1}}\right)\) ，则 \(\mathbf{S}\left( {u,v}\right)\) 在控制点 \({\mathbf{P}}_{i,j},{i}_{0} - p \leq  i \leq  {i}_{0}\) 和 \({j}_{0} - q \leq  j \leq  {j}_{0}\) 的凸包内(见图 3.21)；这由 P3.12、P3.13 和 P3.16 得出；

\begin{center}
\begin{minipage}{18cm}
\vspace*{.6\baselineskip}
\centering
\BookFigure[max width=\linewidth]{ch03/fig_3_23.tex}
\par\smallskip
图 3.21. 三次 \(\times\) 二次示意曲面及其强凸包性质。\\ (a) 虚线围成节点矩形 \([1/3,2/3]^2\) 对应的曲面片；\\ (b) 实线连接其12个活动控制点，外部控制网用虚线表示。
\end{minipage}
\end{center}



P3.23 如果进行三角剖分，控制网形成对曲面的分段平面近似；与曲线的情况一样，度数越低，近似效果越好(见图 3.22a 和 3.22b)；

\begin{center}
\begin{minipage}{18cm}
\vspace*{.6\baselineskip}
\centering
\BookFigure[max width=\linewidth]{ch03/fig_3_24.tex}
\par\smallskip
图 3.22. 同一组代表性 \(5\times5\) 控制点构成的两张曲面。\\ (a) 双二次B样条曲面；(b) 双四次Bézier曲面（\(p=q=4\)）。
\end{minipage}
\end{center}



\begin{center}
\begin{minipage}{18cm}
\vspace*{.6\baselineskip}
\centering
\BookFigure[max width=\linewidth]{ch03/fig_3_25.tex}
\par\smallskip
图 3.23. 单个控制点的局部修改，采用代表性 \(6\times8\) 控制网。\\ (a) 平面二次 \(\times\) 三次曲面，\(U=(0,0,0,1/4,1/2,3/4,1,1,1)\)，\\ \(V=(0,0,0,0,1/5,2/5,3/5,4/5,1,1,1,1)\)。\\ (b) 仅移动 \(\mathbf P_{3,5}\)，影响范围为 \([1/4,1)\times[2/5,1)\)。控制网上移显示以便观察。
\end{minipage}
\end{center}



P3.24 局部修改方案:如果 \({\mathbf{P}}_{i,j}\) 被移动，它仅在矩形 \(\left\lbrack  {{u}_{i},{u}_{i + p + 1}}\right)  \times  \left\lbrack  {{v}_{j},{v}_{j + q + 1}}\right)\) 内影响曲面；这由 P3.15 得出。现在考虑图 3.23a 和 3.23b:初始曲面是平坦的，因为所有控制点都位于同一平面内(P3.22)；控制网偏离曲面以便于可视化。当 \({\mathbf{P}}_{3,5}\) 被移动时，它仅在矩形 \(\left\lbrack  {1/4,1}\right)  \times  \left\lbrack  {2/5,1}\right)\) 内影响曲面形状；

P3.25 \(\mathbf{S}\left( {u,v}\right)\) 的连续性和可微性源于基函数的连续性。特别地， \(\mathbf{S}\left( {u,v}\right)\) 在重数为 \(k\) 的 \(u\left( v\right)\) 节点的 \(u\left( v\right)\) 方向上是 \(p - k\left( {q - k}\right)\) 次可微的。图 3.24 显示了在节点向量 \(U = \{ 0,0,0,1/2,1/2,1,1,1\}\) 和 \(V = \{ 0,0,0,0,1/2,1,1,1,1\}\) 上定义的二次 \(\times\) 三次曲面。注意表面中的折痕，对应于节点线 \(u = 1/2\) 。当然，正如曲线的情况一样，可以以某种方式定位控制点，使其抵消基函数中的不连续性。通过使用重合的控制点，可以在基函数中没有相应不连续性的地方创建视觉不连续性；图 3.25 显示了这样一个表面，它是三次的且没有重节点。因此，二阶偏导数处处连续。折痕是由于多个控制点造成的。

\begin{center}
\begin{minipage}{18cm}
\vspace*{.6\baselineskip}
\centering
\BookFigure[max width=\linewidth]{ch03/fig_3_26.tex}
\par\smallskip
图 3.24. 带折痕的二次 \(\times\) 三次示意曲面。\\ \(U=(0,0,0,1/2,1/2,1,1,1)\)，\(V=(0,0,0,0,1/2,1,1,1,1)\)。\\ \(u=1/2\) 为二重节点，对应图中的折痕。
\end{minipage}
\end{center}

\begin{center}
\begin{minipage}{18cm}
\vspace*{.6\baselineskip}
\centering
\BookFigure[max width=\linewidth]{ch03/fig_3_27.tex}
\par\smallskip
图 3.25. 重合控制点形成折痕的双三次示意曲面。\\ \(U=(0,0,0,0,1/4,1/2,3/4,1,1,1,1)\)，\(V=(0,0,0,0,1/2,1,1,1,1)\)。\\ 对所有 \(0\leq j\leq4\)，\(\mathbf P_{2,j}=\mathbf P_{3,j}=\mathbf P_{4,j}\)；内部节点均为单节点。
\end{minipage}
\end{center}

我们在此指出，目前尚无已知的 B-spline 曲面变差减小性质(见 \cite{Prau92})。

\(\mathbf{S}\left( {u,v}\right)\) 上的等参数曲线\BookIndexMark{p0069}是以类似于 Bézier 曲面的方式获得的。固定 \(u = {u}_{0}\)

\[
{\mathbf{C}}_{{u}_{0}}\left( v\right)  = \mathbf{S}\left( {{u}_{0},v}\right)  = \mathop{\sum }\limits_{{i = 0}}^{n}\mathop{\sum }\limits_{{j = 0}}^{m}{N}_{i,p}\left( {u}_{0}\right) {N}_{j,q}\left( v\right) {\mathbf{P}}_{i,j}
\]

\[
= \mathop{\sum }\limits_{{j = 0}}^{m}{N}_{j,q}\left( v\right) \left( {\mathop{\sum }\limits_{{i = 0}}^{n}{N}_{i,p}\left( {u}_{0}\right) {\mathbf{P}}_{i,j}}\right)  = \mathop{\sum }\limits_{{j = 0}}^{m}{N}_{j,q}\left( v\right) {\mathbf{Q}}_{j}\left( {u}_{0}\right)  \tag{3.14}
\]





其中

\[
{\mathbf{Q}}_{j}\left( {u}_{0}\right)  = \mathop{\sum }\limits_{{i = 0}}^{n}{N}_{i,p}\left( {u}_{0}\right) {\mathbf{P}}_{i,j}
\]

类似地

\[
{\mathbf{C}}_{{v}_{0}}\left( u\right)  = \mathop{\sum }\limits_{{i = 0}}^{n}{N}_{i,p}\left( u\right) {\mathbf{Q}}_{i}\left( {v}_{0}\right)
\]

其中

\[
{\mathbf{Q}}_{i}\left( {v}_{0}\right)  = \mathop{\sum }\limits_{{j = 0}}^{m}{N}_{j,q}\left( {v}_{0}\right) {\mathbf{P}}_{i,j} \tag{3.15}
\]

是一个 \(u\) 等值线在 \(\mathbf{S}\left( {u,v}\right) .{\mathbf{C}}_{{u}_{0}}\left( v\right)\) 上，是一个 \(q\) 次 B 样条曲线在 \(V\) 上，而 \({\mathbf{C}}_{{v}_{0}}\left( u\right)\) 是一个 \(p\) 次 B 样条曲线在 \(U\) 上。点 \(\mathbf{S}\left( {{u}_{0},{v}_{0}}\right)\) 位于 \({\mathbf{C}}_{{u}_{0}}\left( v\right)\) 和 \({\mathbf{C}}_{{v}_{0}}\left( u\right)\) 的交点上。图 3.20-3.25 中显示的所有线条都是等高线。



\section{B样条曲面的导数}\BookIndexMark{p0070}

设 (u, v) 为固定值。一般来说，人们感兴趣的是计算 \(\mathbf{S}\left( {u,v}\right)\) 的所有偏导数，直到包括 \(d\) 阶，即

\[
\frac{{\partial }^{k + l}}{{\partial }^{k}u{\partial }^{l}v}\mathbf{S}\left( {u,v}\right) \;0 \leq  k + l \leq  d \tag{3.16}
\]

对于曲线，我们通过计算基函数的导数来获得这些导数(见方程 [2.9] 和 [2.10] 以及算法 A2.3)。特别地，

\[
\frac{{\partial }^{k + l}}{{\partial }^{k}u{\partial }^{l}v}\mathbf{S}\left( {u,v}\right)  = \mathop{\sum }\limits_{{i = 0}}^{n}\mathop{\sum }\limits_{{j = 0}}^{m}{N}_{i,p}^{\left( k\right) }{N}_{j,q}^{\left( l\right) }{\mathbf{P}}_{i,j} \tag{3.17}
\]

算法 A3.6 计算 B 样条曲面上的点及所有偏导数，直到包括 \(d\) 阶(允许 \(d > p,q\) )。与算法 A3.5 类似，这是一个五步过程，最后一步是形式为向量/矩阵/向量的乘法。

\[
\frac{{\partial }^{k + l}}{{\partial }^{k}u{\partial }^{l}v}\mathbf{S}\left( {u,v}\right)  = {\left\lbrack  {N}_{r,p}^{\left( k\right) }\left( u\right) \right\rbrack  }^{T}\left\lbrack  {\mathbf{P}}_{r,s}\right\rbrack  \left\lbrack  {{N}_{s,q}^{\left( l\right) }\left( v\right) }\right\rbrack
\]

\[
0 \leq  k + l \leq  d\;\text{ uspan } - p \leq  r \leq  u\text{ span }
\]

\[
v\text{ span } - q \leq  s \leq  v\text{ span } \tag{3.18}
\]

输出是数组 SKL[] []，其中 SKL[k] [1] 是 \(\mathbf{S}\left( {u,v}\right)\) 关于 \({uk}\) 和 \({vl}\) 的导数。对于固定的 \(k,0 \leq  k \leq  d\) ，局部数组 temp [] 存储向量/矩阵乘积 \({\left\lbrack  {N}_{r,p}{}^{\left( k\right) }\left( u\right) \right\rbrack  }^{T}\left\lbrack  {\mathbf{P}}_{r,s}\right\rbrack\) ，同时它与 \(\left\lbrack  {{N}_{s,q}{}^{\left( l\right) }\left( v\right) }\right\rbrack\) 相乘，针对 \(0 \leq  l \leq  d - k\) 。数组 \(\mathrm{{Nu}}\left\lbrack  \right\rbrack  \left\lbrack  \right\rbrack\) 和 \(\mathrm{{Nv}}\left\lbrack  \right\rbrack\) [] 用于存储基函数的导数。

% book-code: algorithm A3.6
\par\medskip\noindent\textbf{算法 A3.6（伪代码）}\BookIndexMark{p0071}\par
\begin{minted}[escapeinside=@@,mathescape]{c}
SurfaceDerivsAlg1(n, p, U, m, q, V, P, u, v, d, SKL)
{ /* Compute surface derivatives */
    /* Input: n, p, U, m, q, V, P, u, v, d */
    /* Output: SKL */
    du = min(d, p);
    for ( k = p + 1;k <= d;k ++ )
        for (l=0; l<=d-k; l++) SKL[k][l]= 0.0 ;
    dv = min(d, q);
    for (l=q+1; l<=d; l++)
        for ( k = 0;k <= d - l;k ++ ) SKL[k][l]= 0.0 ;
    uspan = FindSpan(n, p, u, U);
    DersBasisFuns(uspan, u, p, du, U, Nu);
    vspan = FindSpan(m, q, v, V);
    DersBasisFuns(vspan, v, q, dv, V, Nv);
    for ( k = 0 ; k <= du ; k ++ )
    {
        for (s = 0 ; s<=q ; s++)
        {
            temp[s] = 0.0 ;
            for (r=0; r<=p; r++)
                temp[s] = temp[s] + Nu[k][r]*P[uspan-p+r][vspan-q+s];
        }
        dd = min(d-k, dv);
        for (l=0; l<=dd; l++)
        {
            SKL[k][l] = 0.0 ;
            for ( s = 0;s <= q;s ++ )
                SKL[k][l] = SKL[k][l] + Nv[l][s]*temp[s];
        }
    }
}
\end{minted}

图 3.26 显示了一个双立方曲面及其一阶和二阶偏导数。请注意，为了更好地可视化，导数被缩小了 \(1/2\) 。

\begin{center}
\begin{minipage}{18cm}
\vspace*{.6\baselineskip}
\centering
\BookFigure[max width=\linewidth]{ch03/fig_3_28.tex}
\par\smallskip
图 3.26. 双三次示意曲面及其一阶、二阶偏导数。\\ \(U=V=(0,0,0,0,1/2,1,1,1,1)\)，求导位置为 \((u,v)=(3/4,2/5)\)。\\ 两条虚线为相应的等参数曲线；五个导数向量均按 \(1/2\) 缩放显示。
\end{minipage}
\end{center}

让我们正式区分 \(\mathbf{S}\left( {u,v}\right)\) 。关于 \(u\) 我们有

\[
{\mathbf{S}}_{u}\left( {u,v}\right)  = \frac{\partial }{\partial u}\mathbf{S}\left( {u,v}\right)  = \mathop{\sum }\limits_{{j = 0}}^{m}{N}_{j,q}\left( v\right) \left( {\frac{\partial }{\partial u}\mathop{\sum }\limits_{{i = 0}}^{n}{N}_{i,p}\left( u\right) {\mathbf{P}}_{i,j}}\right)  \tag{3.19}
\]

\[
= \mathop{\sum }\limits_{{j = 0}}^{m}{N}_{j,q}\left( v\right) \left( {\frac{\partial }{\partial u}{\mathbf{C}}_{j}\left( u\right) }\right)
\]

其中

\[
{\mathbf{C}}_{j}\left( u\right)  = \mathop{\sum }\limits_{{i = 0}}^{n}{N}_{i,p}\left( u\right) {\mathbf{P}}_{i,j}\;j = 0,\ldots ,m
\]



是 B 样条曲线。将公式 (3.6) 应用到每个 \({\mathbf{C}}_{j}\left( u\right)\) 并进行替换

代入公式 (3.19)，我们得到

\[
{\mathbf{S}}_{u}\left( {u,v}\right)  = \mathop{\sum }\limits_{{i = 0}}^{{n - 1}}\mathop{\sum }\limits_{{j = 0}}^{m}{N}_{i,p - 1}\left( u\right) {N}_{j,q}\left( v\right) {\mathbf{P}}_{i,j}^{\left( 1,0\right) } \tag{3.20}
\]

其中

\[
{\mathbf{P}}_{i,j}^{\left( 1,0\right) } = p\frac{{\mathbf{P}}_{i + 1,j} - {\mathbf{P}}_{i,j}}{{u}_{i + p + 1} - {u}_{i + 1}}\;\text{ (see Eq. [3.4]) }
\]

\[
{U}^{\left( 1\right) } = \left\{  {\underset{p}{\underbrace{0,\ldots ,0}},{u}_{p + 1},\ldots ,{u}_{r - p - 1},\underset{p}{\underbrace{1,\ldots ,1}}}\right\}
\]

\[
{V}^{\left( 0\right) } = V
\]

类似地 \(\;{\mathbf{S}}_{v}\left( {u,v}\right)  = \mathop{\sum }\limits_{{i = 0}}^{n}\mathop{\sum }\limits_{{j = 0}}^{{m - 1}}{N}_{i,p}\left( u\right) {N}_{j,q - 1}\left( v\right) {\mathbf{P}}_{i,j}^{\left( 0,1\right) }\)(3.21)

其中

\[
{\mathbf{P}}_{i,j}^{\left( 0,1\right) } = q\frac{{\mathbf{P}}_{i,j + 1} - {\mathbf{P}}_{i,j}}{{v}_{j + q + 1} - {v}_{j + 1}}
\]

\[
{U}^{\left( 0\right) } = U
\]

\[
{V}^{\left( 1\right) } = \left\{  {\underset{q}{\underbrace{0,\ldots ,0}},{v}_{q + 1},\ldots ,{v}_{s - q - 1},\underset{q}{\underbrace{1,\ldots ,1}}}\right\}
\]

首先应用公式 (3.20)，然后应用公式 (3.21) 得到

\[
{\mathbf{S}}_{uv}\left( {u,v}\right)  = \mathop{\sum }\limits_{{i = 0}}^{{n - 1}}\mathop{\sum }\limits_{{j = 0}}^{{m - 1}}{N}_{i,p - 1}\left( u\right) {N}_{j,q - 1}\left( v\right) {\mathbf{P}}_{i,j}^{\left( 1,1\right) } \tag{3.22}
\]

其中

\[
{\mathbf{P}}_{i,j}^{\left( 1,1\right) } = q\frac{{\mathbf{P}}_{i,j + 1}^{\left( 1,0\right) } - {\mathbf{P}}_{i,j}^{\left( 1,0\right) }}{{v}_{j + q + 1} - {v}_{j + 1}}
\]

并且 \({U}^{\left( 1\right) }\) 和 \({V}^{\left( 1\right) }\) 的定义如前所述。

一般来说

\[
\frac{{\partial }^{k + l}}{{\partial }^{k}u{\partial }^{l}v}\mathbf{S}\left( {u,v}\right)  = \mathop{\sum }\limits_{{i = 0}}^{{n - k}}\mathop{\sum }\limits_{{j = 0}}^{{m - l}}{N}_{i,p - k}\left( u\right) {N}_{j,q - l}\left( v\right) {\mathbf{P}}_{i,j}^{\left( k,l\right) } \tag{3.23}
\]

在哪里

\[
{\mathbf{P}}_{i,j}^{\left( k,l\right) } = \left( {q - l + 1}\right) \frac{{\mathbf{P}}_{i,j + 1}^{\left( k,l - 1\right) } - {\mathbf{P}}_{i,j}^{\left( k,l - 1\right) }}{{v}_{j + q + 1} - {v}_{j + l}}
\]

114 B样条曲线和曲面

使用公式 (3.20)-(3.23)，我们推导出角导数的有用公式。例如，在角点 \(\left( {u,v}\right)  = \left( {0,0}\right)\) 处，我们有

\[
{\mathbf{S}}_{u}\left( {0,0}\right)  = {\mathbf{P}}_{0,0}^{\left( 1,0\right) } = \frac{p}{{u}_{p + 1}}\left( {{\mathbf{P}}_{1,0} - {\mathbf{P}}_{0,0}}\right)
\]

\[
{\mathbf{S}}_{v}\left( {0,0}\right)  = {\mathbf{P}}_{0,0}^{\left( 0,1\right) } = \frac{q}{{v}_{q + 1}}\left( {{\mathbf{P}}_{0,1} - {\mathbf{P}}_{0,0}}\right)  \tag{3.24}
\]

\[
{\mathbf{S}}_{uv}\left( {0,0}\right)  = {\mathbf{P}}_{0,0}^{\left( 1,1\right) } = \frac{q}{{v}_{q + 1}}\left( {{\mathbf{P}}_{0,1}^{\left( 1,0\right) } - {\mathbf{P}}_{0,0}^{\left( 1,0\right) }}\right)
\]

\[
= \frac{pq}{{u}_{p + 1}{v}_{q + 1}}\left( {{\mathbf{P}}_{1,1} - {\mathbf{P}}_{0,1} - {\mathbf{P}}_{1,0} + {\mathbf{P}}_{0,0}}\right)
\]

现在设 \({u}_{0} = 0\) 和 \({v}_{0} = 0\) 。根据基函数的性质，很容易看出等曲线 \({\mathbf{C}}_{{u}_{0}}\left( v\right)\) 和 \({\mathbf{C}}_{{v}_{0}}\left( u\right)\) 是由以下公式给出的

\[
{\mathbf{C}}_{{u}_{0}}\left( v\right)  = \mathop{\sum }\limits_{{j = 0}}^{m}{N}_{j,q}\left( v\right) {\mathbf{P}}_{0,j}\;{\mathbf{C}}_{{v}_{0}}\left( u\right)  = \mathop{\sum }\limits_{{i = 0}}^{n}{N}_{i,p}\left( u\right) {\mathbf{P}}_{i,0}
\]

从公式 (3.7) 可以得出

\[
{\mathbf{S}}_{u}\left( {0,0}\right)  = {\mathbf{C}}_{{v}_{0}}^{\prime }\left( 0\right) \;{\mathbf{S}}_{v}\left( {0,0}\right)  = {\mathbf{C}}_{{u}_{0}}^{\prime }\left( 0\right)
\]

算法 A3.7 计算所有(或可选的一些)控制点 \({\mathbf{P}}_{i,j}^{\left( k,l\right) }\) ，导数曲面最高到 \(d\left( {0 \leq  k + l \leq  d}\right)\) 阶。该算法基于公式 (3.23) 和算法 A3.3。输出为数组 PKL[][][][]，其中 \(\operatorname{PKL}\left\lbrack  \mathrm{k}\right\rbrack  \left\lbrack  1\right\rbrack  \left\lbrack  i\right\rbrack  \left\lbrack  j\right\rbrack\) 是曲面的第 \(i,j\) 个控制点，相对于 \(u\) 导数求导 \(k\) 次，相对于 \(l\) 导数求导 \(v\) 次。

% book-code: algorithm A3.7
\par\medskip\noindent\textbf{算法 A3.7（伪代码）}\BookIndexMark{p0072}\par
\begin{minted}[escapeinside=@@,mathescape]{c}
SurfaceDerivCpts( n,p,U,m,q,V,P,d,r1,r2,s1,s2,PKL)
{ /* Compute control points of derivative surfaces */
    /* Input: n, p, U, m, q, V, P, d, r1, r2, s1, s2 */
    /* Output: PKL */
    @${du} = \min \left( {d,p}\right) ;\;{dv} = \min \left( {d,q}\right) ;$@
    r = r2-r1; s = s2-s1;
    for (j=s1; j<=s2; j++)
    {
        CurveDerivCpts(n, p, U,&P[][j], du, r1, r2, temp);
        for (k=0; k<=du; k++)
            for (i = 0; i<=r-k; i++)
                PKL[k][0][i][j-s1] = temp[k][i];
    }
    for ( k = 0;k < du;k ++ )
        for (i = 0; i<=r-k; i++)
        {
            dd = min(d-k, dv);
            CurveDerivCpts(m, q,&V[s1],&PKL[k][0][i][], dd,0, s, temp);
            for (l=1; l<=dd; l++)
                for (j=0; j<=s-1; j++)
                    PKL[k][l][i][j] = temp[l][j];
        }
}
\end{minted}

算法 A3.8 计算 B样条曲面上的点及所有偏导数，最高到并包括 \(d\) 阶，固定参数 (u, v)(与算法 A3.6 比较)。允许使用 \(d > p,q\) 。输出时，SKL[k][1] 是相对于 \(u\) 导数求导 \(\mathbf{S}\left( {u,v}\right) k\) 次，相对于 \(v\) 导数求导 \(l\) 次。

% book-code: algorithm A3.8
\par\medskip\noindent\textbf{算法 A3.8（伪代码）}\BookIndexMark{p0073}\par
\begin{minted}[escapeinside=@@,mathescape]{c}
SurfaceDerivsAlg2(n, p, U, m, q, V, P, u, v, d, SKL)
{ /* Compute surface derivatives */
    /* Input: n, p, U, m, q, V, P, u, v, d */
    /* Output: SKL */
    du = min(d, p);
    for (k=p+1; k<=d; k++)
        for (l=0; l<=d-k; l++) SKL[k][l] = 0.0;
    dv = min(d, q);
    for (l=q+1; l<=d; l++)
        for ( k = 0;k <= d - l;k ++ ) SKL[k][l]= 0.0 ;
    uspan = FindSpan(n, p, u, U);
    AllBasisFuns(uspan, u, p, U, Nu);
    vspan = FindSpan(m, q, v, V);
    AllBasisFuns(vspan, v, q, V, Nv);
    SurfaceDerivCpts (n,p,U,m,q,V,P,d,uspan - p,uspan,
        vspan-q, vspan, PKL);
    for ( k = 0 ; k <= du ; k ++ )
    {
        dd = min(d-k, dv);
        for (l=0; l<=dd; l++)
        {
            SKL[k][l] = 0.0 ;
            for (i = 0; i<=q-l; i++)
            {
                tmp = 0.0 ;
                for (j=0; j<=p-k; j++)
                    tmp = tmp + Nu[j][p-k]*PKL[k][l][j][i];
                SKL[k][l] = SKL[k][l] + Nv[i][q-l]*tmp;
            }
        }
    }
}
\end{minted}

\section*{练习}

\section*{3.1. 为什么二次曲线在节点处接触其控制多边形？}

3.2. 如果一条二次曲线有一个拐点，它必须位于一个节点上(见图)。为什么？

3.3. 构造一个具有尖点的 \({\mathbf{C}}^{2}\) 连续三次曲线。

3.4. 设三次曲线由 \(\mathbf{C}\left( u\right)  = \mathop{\sum }\limits_{{i = 0}}^{7}{N}_{i,3}\left( u\right) {\mathbf{P}}_{i}\) 和节点向量 \(U = \{ 0,0,0,0,1/4,1/4,2/3,3/4,1,1,1,1\} .\) 定义。

a. 假设 \({\mathbf{P}}_{i}\) 的一些任意位置，并绘制曲线。

b. 点 \(\mathbf{C}\left( {1/4}\right)\) 在哪里？

c. 如果 \({\mathbf{P}}_{2}\) 被移动， \(\mathbf{C}\left( u\right)\) 在 \(\left\lbrack  {0,1}\right\rbrack\) 的哪个子区间受到影响？如果 \({\mathbf{P}}_{5}\) 被移动，哪个子区间受到影响？

d. 哪些控制点在区间 \(u \in  \left\lbrack  {1/4,2/3}\right)\) 上影响曲线形状？在区间 \(u \in  \left\lbrack  {2/3,3/4}\right)\) 上呢？

3.5. 设 \(\mathbf{C}\left( u\right)  = \mathop{\sum }\limits_{{i = 0}}^{3}{N}_{i,2}\left( u\right) {\mathbf{P}}_{i}\) ，其中 \(U = \{ 0,0,0,1/2,1,1,1\}\) 和 \({\mathbf{P}}_{0} = \left( {-1,0}\right)\) ， \({\mathbf{P}}_{1} = \left( {-1,1}\right) ,{\mathbf{P}}_{2} = \left( {1,1}\right)\) 和 \({\mathbf{P}}_{3} = \left( {1,0}\right)\) 。绘制 \(\mathbf{C}\left( u\right)\) 。计算 \({\mathbf{C}}^{\prime }\left( u\right)\) ，即其控制点和节点向量。绘制 \({\mathbf{C}}^{\prime }\left( u\right)\) 。

3.6. 对于图 3.16a 的三次曲线，假设控制点为 \(\left\{  {\mathbf{P}}_{i}\right\}   = \{ \left( {0,0}\right)\) ， \(\left( {1,2}\right) ,\left( {3,4}\right) ,\left( {5,2}\right) ,\left( {5,0}\right) ,\left( {8,0}\right) ,\left( {9,3}\right) \}\) 。绘制曲线。使用公式 (3.8) 计算二阶导数曲线 \({\mathbf{C}}^{\left( 2\right) }\left( u\right)\) 。绘制 \({\mathbf{C}}^{\left( 2\right) }\left( u\right)\) 。设 \({\mathbf{P}}_{0}^{\left( 2\right) }\) 和 \({\mathbf{P}}_{4}^{\left( 2\right) }\) 为 \({\mathbf{C}}^{\left( 2\right) }\left( u\right)\) 的第一个和最后一个控制点，即 \({\mathbf{P}}_{0}^{\left( 2\right) } = {\mathbf{C}}^{\left( 2\right) }\left( 0\right)\) 和 \({\mathbf{P}}_{4}^{\left( 2\right) } = {\mathbf{C}}^{\left( 2\right) }\left( 1\right)\) 。使用公式 (3.9) 和 (3.10) 计算 \({\mathbf{C}}^{\left( 2\right) }\left( 0\right)\) 和 \({\mathbf{C}}^{\left( 2\right) }\left( 1\right)\) 。

3.7. 可以使用多个节点或多个控制点在表面上创建折痕(见图 3.24 和 3.25)。如果您想要一个折痕长度小于表面全长，您会使用哪种方法？构造并绘制这样的示例。

3.8. 考虑 B-spline 表面 \(\mathbf{S}\left( {u,v}\right)  = \mathop{\sum }\limits_{{i = 0}}^{3}\mathop{\sum }\limits_{{j = 0}}^{2}{N}_{i,2}\left( u\right) {N}_{j,2}\left( v\right) {\mathbf{P}}_{i,j}\)

其中

\[
U = \{ 0,0,0,1/2,1,1,1\}
\]

\[
V = \{ 0,0,0,1,1,1\}
\]

和

\[
{\mathbf{P}}_{0,0} = \left( {0,0,0}\right) \;{\mathbf{P}}_{1,0} = \left( {3,0,3}\right) \;{\mathbf{P}}_{2,0} = \left( {6,0,3}\right) \;{\mathbf{P}}_{3,0} = \left( {9,0,0}\right)
\]

\[
{\mathbf{P}}_{0,1} = \left( {0,2,2}\right) \;{\mathbf{P}}_{1,1} = \left( {3,2,5}\right) \;{\mathbf{P}}_{2,1} = \left( {6,2,5}\right) \;{\mathbf{P}}_{3,1} = \left( {9,2,2}\right)
\]

\[
{\mathbf{P}}_{0,2} = \left( {0,4,0}\right) \;{\mathbf{P}}_{1,2} = \left( {3,4,3}\right) \;{\mathbf{P}}_{2,2} = \left( {6,4,3}\right) \;{\mathbf{P}}_{3,2} = \left( {9,4,0}\right)
\]

通过评估非零 \(\mathbf{B}\) -样条基函数并将其与适当的控制点相乘来计算 \(\mathbf{S}\left( {3/{10},6/{10}}\right)\) 。

3.9. 推导 \({\mathbf{S}}_{u}\left( {u,v}\right) ,{\mathbf{S}}_{v}\left( {u,v}\right)\) 和 \({\mathbf{S}}_{uv}\left( {u,v}\right)\) 在三个角 \(\left( {u,v}\right)  = \left( {0,1}\right) ,\left( {1,0}\right)\) 和 (1,1) 的表达式(见公式 [3.24])。

3.10. 设 \(\mathbf{S}\left( {u,v}\right)\) 如习题 3.8 所示。绘制该曲面。使用公式 (3.20) 和 (3.21)，计算曲面 \({\mathbf{S}}_{u}\left( {u,v}\right)\) 和 \({\mathbf{S}}_{v}\left( {u,v}\right)\) 。绘制这两个曲面。使用公式 (3.22) 和在习题 3.9 中推导的表达式，计算曲面四个角的混合偏导数 \({\mathbf{S}}_{uv}\left( {u,v}\right)\) 。这四个向量的几何意义是什么？




\clearpage
\end{document}
